\documentclass[10pt]{article}
\usepackage[a4paper,margin=0.6in]{geometry}
\usepackage{amsmath,amssymb,amsfonts,mathtools}
\usepackage[dvipsnames]{xcolor}
\usepackage{array,booktabs,float,longtable}
\usepackage[colorlinks=true,linkcolor=blue,citecolor=blue]{hyperref}
\allowdisplaybreaks
\numberwithin{equation}{section}
\setlength{\jot}{3pt}
\setlength{\emergencystretch}{3em}
\newcommand{\new}[1]{{\color{blue}#1}}
\newcommand{\old}[1]{{\color{red}#1}}
\newcommand{\NN}{\mathcal N}
\newcommand{\ok}{\textcolor{ForestGreen}{correct}}

\title{Check of \texttt{Evolution\_terms.tex}:\\ the evolution terms, JIMWLK, BFKL and their colour-singlet forms}
\date{}

\begin{document}
\maketitle
\tableofcontents

%==========================================================================================
\section{Summary}\label{sec:summary}
Throughout, $\NN=\frac{1}{(2\pi)^3}\frac{g^4}{16\pi^5}\frac{1}{k^+}$, the prefactor of every equation of your file. JIMWLK$\otimes$LO and BFKL$\otimes$LO are the targets of the earlier notes (Eqs.~(3.1), (3.2) of \texttt{Evolution\_singlet.pdf}).

\paragraph{What is correct.}
\begin{enumerate}
\item \textbf{Sec.~1, genuine terms.} All 8 equations agree with the checked genuine terms (Group I Rows II, V; Group II Rows I, IV; Group III Row III), Table~\ref{tab:gen}.
\item \textbf{Sec.~2, from three and four $\rho$'s.} All 15 equations are identical, term by term, to the lists of the previous version, which were checked against the reordered three- and four-$\rho$ terms (Table~\ref{tab:sec2}).
\item \textbf{Sec.~3, JIMWLK.} Your collected form, Eq.~(3.3) of your file, \emph{is} JIMWLK$\otimes$LO, exactly, for any adjoint $U$ ($3.4\times10^{-17}$). $dN_1+dN_2$ equal it with your convention $U^{bc}=U^{cb}$.
\item \textbf{Sec.~4, BFKL.} Your collected form, Eq.~(4.2), \emph{is} BFKL$\otimes$LO, exactly, for any $U$ ($8.7\times10^{-16}$). Eq.~(4.1) equals it with $U^{bc}=U^{cb}$.
\item \textbf{Colour-singlet forms} (Eqs.~(3.4), (4.3)). Both are correct for any $U$:
\begin{itemize}
\item they equal the singlet projection of your own collected forms ($2.3\times10^{-15}$, $6.0\times10^{-16}$);
\item they equal Eqs.~(3.1), (3.2) of \texttt{Evolution\_singlet.pdf} ($1.4\times10^{-15}$, $4.2\times10^{-16}$).
\end{itemize}
\end{enumerate}

\paragraph{What is not.}
\begin{enumerate}
\item[\textbf{P1}] \textbf{Secs.~1 and 2 do not add up to Sec.~3} (this is M1 of \texttt{Final\_evolution\_check.pdf}, still open). The ``From $4\,\rho$'s'' list of Sec.~2 is the one of the previous version. The two-$\rho$ lists of \texttt{4rho.tex} (30 terms) are still missing.
\begin{itemize}
\item Without them, Secs.~1+2 differ from JIMWLK$\otimes$LO by 8--54\% (relative, at random points, your convention), and by 27--170\% in the singlet projection.
\item The differing kernel structures are $\mathcal K(x-w)\cdot\mathcal K(x'-w')$ times $\frac{1}{(x-z)^2}$, $\frac{1}{(x'-z)^2}$ and $K(x,x',z)$.
\item With them, Secs.~1+2 equal Sec.~3 exactly ($8.4\times10^{-17}$). The 30 terms are written out in Eq.~\eqref{eq:four}.
\end{itemize}
\item[\textbf{P2}] \textbf{The BFKL section carries the wrong logarithm.}
\begin{itemize}
\item As written, both Sec.~3 and Sec.~4 are $d^3N_{\rm NLO}/d^3k\propto\log\frac{\vee}{\Lambda}$. Added, they would give $\log\frac{\vee}{\Lambda}\,(\mathcal J+\mathcal B)$.
\item But the coefficient of $\log\frac{\vee}{\Lambda}$ is $\mathcal J$ alone: it is the sum of Secs.~1, 2, which is Sec.~3 (P1 fixed).
\item $\mathcal B$ is the coefficient of $\log\vee$. It comes from the upper limits of the $p^+$ integrals and multiplies $\log\frac{\vee}{k^+}$ (Sec.~\ref{sec:p2}).
\item In Sec.~4, \old{$\log\frac{\vee}{\Lambda}$} $\to$ \new{$\log\frac{\vee}{k^+}$}, and the total is Eq.~\eqref{eq:total}.
\end{itemize}
\item[\textbf{P3}] \textbf{Normalization of the singlet forms.} The bracket is right but the prefactor is not complete (Sec.~\ref{sec:p3}):
\begin{itemize}
\item $\frac{1}{N_c^2-1}\langle\rho^f(x)\rho^f(y)\rangle$ is missing;
\item the integral is written over $x'$, the integrand uses $y$;
\item $K(a,b,z)$ is used before it is defined, and only $K(y,y,z)$ and $K(x,y,z)$ are defined.
\end{itemize}
\item[\textbf{P4}] \textbf{Labels} (Sec.~\ref{sec:p4}).
\begin{itemize}
\item In Group I Row V, the fourth equation is labelled $d^3N_{2b\,\rm finite}$ a second time. It is the counterterm: it equals $-d^3N_{2b\,\rm UV}$ term by term, so the two cancel ($4.9\times10^{-17}$).
\item ``(with c.c.)'' in the headings means that the c.c.\ is already written out. Adding it once more would be wrong: for Group III Row III it would change the result by $8.3\times10^{-2}$.
\end{itemize}
\end{enumerate}

\paragraph{The convention $U^{bc}=U^{cb}$.} Some of your forms are correct only with your convention $U^{bc}=U^{cb}$ ($U^\dagger=U$), as found in \texttt{Evolution\_singlet.pdf}:
\begin{itemize}
\item the Sec.~2 lists;
\item the genuine Group II Row IV and Group III Row III;
\item $dN_1+dN_2$ of Sec.~3;
\item Eq.~(4.1).
\end{itemize}
The collected forms (3.3), (4.2) and the singlet forms (3.4), (4.3) are correct for any $U$. Nothing has to be changed for this, as long as the convention is stated.

\paragraph{Corrected copy.} \texttt{Evolution\_terms\_corrected.tex} (and \texttt{.pdf}) is your file with P1--P4 fixed, removed text in \old{red}, new text in \new{blue}. Parsed back, its Secs.~1+2 equal its Sec.~3 and JIMWLK$\otimes$LO ($8.4\times10^{-17}$).

%==========================================================================================
\section{How the check was done}
\begin{itemize}
\item Every term of the file was read by a parser: 624 two-$\rho$ terms in 30 equations, and the 44 + 8 trace monomials of the two singlet forms. The factored forms ``kernel $\times\,\big[\dots\big]$'' were expanded.
\item Each term was converted to the format of the earlier checks (Wilson lines, structure constants, charges, kernels).
\item Two sums of terms are compared at random transverse points $x, x', w, w', z$, with random classical charges $\rho^a(x)$, $\rho^a(x')$ and random SU(3) adjoint Wilson lines. Charges are not projected, so this is stronger than the singlet comparison.
\begin{itemize}
\item The number given is $|A-B|$ divided by the sum of the moduli of all single terms, maximized over three configurations.
\item Agreement means $\lesssim10^{-15}$.
\item ``$U=U^T$'' is your convention $U^{bc}=U^{cb}$; ``any $U$'' is a general adjoint $U$.
\end{itemize}
\item Singlet forms: the traces are evaluated with the same random $U$, with $x\to p_1$, $y\to p_2$, symmetrized.
\end{itemize}

%==========================================================================================
\section{Equation by equation}
\begin{table}[H]\centering\small
\caption{Sec.~1, genuine terms, against the checked genuine terms of \texttt{Final\_evolution\_check.pdf}.}\label{tab:gen}
\begin{tabular}{@{}lccl@{}}\toprule
equation & $U=U^T$ & any $U$ & \\\midrule
(1.1) Group I, Row II & $3.0\times10^{-17}$ & $2.1\times10^{-17}$ & \textcolor{ForestGreen}{agrees}, any $U$\\
(1.2) Group I, Row V, $d^3N_{2a}$ & $6.0\times10^{-17}$ & $7.9\times10^{-17}$ & \textcolor{ForestGreen}{agrees}, any $U$\\
(1.3) Group I, Row V, $d^3N_{2b\,\rm UV}$ & $3.6\times10^{-17}$ & $2.1\times10^{-17}$ & \textcolor{ForestGreen}{agrees}, any $U$\\
(1.4) Group I, Row V, $d^3N_{2b\,\rm finite}$ & $2.1\times10^{-17}$ & $1.3\times10^{-16}$ & \textcolor{ForestGreen}{agrees}, any $U$\\
(1.5) Group I, Row V, second ``$d^3N_{2b\,\rm finite}$'' vs counterterm & $3.6\times10^{-17}$ & $2.1\times10^{-17}$ & \textcolor{ForestGreen}{agrees}, any $U$\\
(1.3) + (1.5), relative to (1.3) (vanishes) & $4.9\times10^{-17}$ & $5.9\times10^{-17}$ & \textcolor{ForestGreen}{agrees}, any $U$\\
(1.2)--(1.5) together vs Row V & $2.5\times10^{-17}$ & $3.1\times10^{-17}$ & \textcolor{ForestGreen}{agrees}, any $U$\\
(1.6) Group II, Row I & $3.4\times10^{-17}$ & $9.7\times10^{-18}$ & \textcolor{ForestGreen}{agrees}, any $U$\\
(1.7) Group II, Row IV & $2.0\times10^{-16}$ & $3.0\times10^{-1}$ & \textcolor{ForestGreen}{agrees} for $U^{bc}=U^{cb}$\\
(1.8) Group III, Row III, second part (as written) & $1.3\times10^{-16}$ & $2.1\times10^{-1}$ & \textcolor{ForestGreen}{agrees} for $U^{bc}=U^{cb}$\\
(1.8) with its c.c.\ added once more & $8.3\times10^{-2}$ & $3.1\times10^{-1}$ & \textcolor{red}{differs}\\
\midrule
all of Sec.~1 & $7.1\times10^{-17}$ & $5.3\times10^{-2}$ & \textcolor{ForestGreen}{agrees} for $U^{bc}=U^{cb}$\\
\bottomrule
\end{tabular}
\end{table}

\begin{table}[H]\centering\small
\caption{Sec.~2 against the same lists of the previous version, which were checked against the reordered three- and four-$\rho$ terms.}\label{tab:sec2}
\begin{tabular}{@{}lcccc@{}}\toprule
equation & terms (new / previous) & $U=U^T$ & any $U$ & \\\midrule
(2.1) From four $\rho$'s & 20 / 20 & $0$ & $0$ & \textcolor{ForestGreen}{identical}\\
(2.2) Group I, Row II & 22 / 22 & $0$ & $0$ & \textcolor{ForestGreen}{identical}\\
(2.3) Group I, Row III, $N_1$ & 21 / 21 & $0$ & $0$ & \textcolor{ForestGreen}{identical}\\
(2.4) Group I, Row III, $N_2$ & 24 / 24 & $0$ & $0$ & \textcolor{ForestGreen}{identical}\\
(2.5) Group I, Row IV, $N_1$ & 33 / 33 & $0$ & $0$ & \textcolor{ForestGreen}{identical}\\
(2.6) Group I, Row IV, $N_2$ & 13 / 13 & $0$ & $0$ & \textcolor{ForestGreen}{identical}\\
(2.7) Group I, Row V & 16 / 16 & $0$ & $0$ & \textcolor{ForestGreen}{identical}\\
(2.8) Group II, Row I & 31 / 31 & $0$ & $0$ & \textcolor{ForestGreen}{identical}\\
(2.9) Group III, Row I, $N_1$ & 30 / 30 & $0$ & $0$ & \textcolor{ForestGreen}{identical}\\
(2.10) Group III, Row I, $N_2$ & 16 / 16 & $0$ & $0$ & \textcolor{ForestGreen}{identical}\\
(2.11) Group III, Row II, $N_1$ & 34 / 34 & $0$ & $0$ & \textcolor{ForestGreen}{identical}\\
(2.12) Group III, Row II, $N_2$ & 11 / 11 & $0$ & $0$ & \textcolor{ForestGreen}{identical}\\
(2.13) Group III, Row III & 16 / 16 & $0$ & $0$ & \textcolor{ForestGreen}{identical}\\
(2.14) Group III, Row V & 24 / 24 & $0$ & $0$ & \textcolor{ForestGreen}{identical}\\
(2.15) Group III, Row VI & 24 / 24 & $0$ & $0$ & \textcolor{ForestGreen}{identical}\\
\bottomrule
\end{tabular}
\end{table}

\begin{table}[H]\centering\small
\caption{Secs.~3, 4 and the totals (normalized by the sum of the moduli of the single terms, Sec.~2).}\label{tab:s34}
\begin{tabular}{@{}lccl@{}}\toprule
comparison & $U=U^T$ & any $U$ & \\\midrule
Secs.~1+2 vs JIMWLK$\otimes$LO & $2.4\times10^{-2}$ & $1.8\times10^{-2}$ & \textcolor{red}{differs}\\
Secs.~1+2 + \texttt{4rho.tex} lists (P1) vs JIMWLK$\otimes$LO & $8.4\times10^{-17}$ & $5.2\times10^{-2}$ & \ok{} for $U^{bc}=U^{cb}$\\
\midrule
(3.1) + (3.2) vs (3.3) & $7.7\times10^{-17}$ & $1.1\times10^{-1}$ & \textcolor{ForestGreen}{agrees} for $U^{bc}=U^{cb}$\\
(3.3) vs JIMWLK$\otimes$LO & $1.5\times10^{-16}$ & $3.4\times10^{-17}$ & \ok{}, any $U$\\
(3.3) vs Secs.~1+2 as written & $2.5\times10^{-2}$ & $1.8\times10^{-2}$ & \textcolor{red}{differs}\\
\midrule
(4.1) vs (4.2) & $4.6\times10^{-17}$ & $1.3\times10^{-1}$ & \textcolor{ForestGreen}{agrees} for $U^{bc}=U^{cb}$\\
(4.2) vs BFKL$\otimes$LO & $2.2\times10^{-15}$ & $8.7\times10^{-16}$ & \ok{}, any $U$\\
(4.1) vs BFKL$\otimes$LO & $2.2\times10^{-15}$ & $9.2\times10^{-2}$ & \ok{} for $U^{bc}=U^{cb}$\\
\bottomrule
\end{tabular}
\end{table}

\begin{table}[H]\centering\small
\caption{The colour-singlet forms (relative difference of the values).}\label{tab:sing}
\begin{tabular}{@{}lccl@{}}\toprule
comparison & $U=U^T$ & any $U$ & \\\midrule
(3.4) vs singlet projection of (3.3) & $1.0\times10^{-15}$ & $2.3\times10^{-15}$ & \textcolor{ForestGreen}{agrees}, any $U$\\
(3.4) vs Eq.~(3.1) of \texttt{Evolution\_singlet.pdf} & $1.6\times10^{-15}$ & $1.4\times10^{-15}$ & \ok{}, any $U$\\
(4.3) vs singlet projection of (4.2) & $1.1\times10^{-15}$ & $6.0\times10^{-16}$ & \textcolor{ForestGreen}{agrees}, any $U$\\
(4.3) vs Eq.~(3.2) of \texttt{Evolution\_singlet.pdf} & $9.3\times10^{-16}$ & $4.2\times10^{-16}$ & \ok{}, any $U$\\
\bottomrule
\end{tabular}
\end{table}

%==========================================================================================
\section{P1: the missing two-\texorpdfstring{$\rho$}{rho} terms of \texttt{4rho.tex}}\label{sec:p1}
These are the two-$\rho$ terms of \texttt{4rho.tex} (``Two $\rho$ contribution'', $dN^{(1)}$, $dN^{(2)}$ with c.c.), with the fixes F1, F2 of \texttt{Final\_evolution\_check.pdf}:
\begin{itemize}
\item \textbf{F1:} in the first part, drop \old{$-8N_c\,\mathcal K(x'-w')\cdot\mathcal K(x'-w)\,K(y,x',z)$};
\item \textbf{F2:} in the second part, $dN^{(1)}$, term 2 has one index four times. The correct term is
\begin{equation*}\new{3f^{abc}f^{deg}\rho^{b'}(x')\rho^c(y)U^{ag}(y)U^{db'}(w')U^{eb}(z)}.\end{equation*}
\end{itemize}
They belong to ``From $4\,\rho$'s'' of Sec.~2. Collected, in your notation (\new{to be added}):
\begin{align}
\frac{d^3N_{\rm NLO}}{d^3k}\bigg|_{\texttt{4rho.tex}}={}&\NN\log\Big[\frac{\vee}{\Lambda}\Big]\int_{x,x',z,w,w'}e^{-ik\cdot(w'-w)}\Bigg[\notag\\
\notag\\
&\quad{\color{blue}+2\,N_c\;\mathcal{K}(x-w)\cdot\mathcal{K}(x'-w')\,\tfrac{1}{(x-z)^{2}}\;U^{ab}(x)\,U^{ac}(w')\,\rho^{c}(x')\,\rho^{b}(x)}\notag\\
&\quad{\color{blue}-2\,N_c\;\mathcal{K}(x-w)\cdot\mathcal{K}(x'-w')\,\tfrac{1}{(x-z)^{2}}\;U^{ab}(x)\,U^{ac}(x')\,\rho^{c}(x')\,\rho^{b}(x)}\notag\\
&\quad{\color{blue}-3\;\mathcal{K}(x-w)\cdot\mathcal{K}(x'-w')\,\tfrac{1}{(x-z)^{2}}\;f^{abc}\,f^{cde}\,U^{ah}(z)\,U^{bl}(w)\,U^{dh}(x)\,U^{ml}(w')\,\rho^{m}(x')\,\rho^{e}(x)}\notag\\
&\quad{\color{blue}+2\;\mathcal{K}(x-w)\cdot\mathcal{K}(x'-w')\,\tfrac{1}{(x-z)^{2}}\;f^{abc}\,f^{cde}\,U^{ah}(z)\,U^{bl}(x)\,U^{dh}(x)\,U^{ml}(w')\,\rho^{m}(x')\,\rho^{e}(x)}\notag\\
&\quad{\color{blue}-3\;\mathcal{K}(x-w)\cdot\mathcal{K}(x'-w')\,\tfrac{1}{(x-z)^{2}}\;f^{abc}\,f^{cde}\,U^{ah}(w)\,U^{bl}(x)\,U^{dl}(z)\,U^{mh}(w')\,\rho^{m}(x')\,\rho^{e}(x)}\notag\\
&\quad{\color{blue}+3\;\mathcal{K}(x-w)\cdot\mathcal{K}(x'-w')\,\tfrac{1}{(x-z)^{2}}\;f^{abc}\,f^{cde}\,U^{ah}(z)\,U^{dl}(w)\,U^{lm}(w')\,U^{hb}(x)\,\rho^{m}(x')\,\rho^{e}(x)}\notag\\
&\quad{\color{blue}+\;\mathcal{K}(x-w)\cdot\mathcal{K}(x'-w')\,\tfrac{1}{(x-z)^{2}}\;f^{abc}\,f^{cde}\,U^{hd}(w)\,U^{hl}(x')\,U^{ma}(z)\,U^{mb}(x)\,\rho^{l}(x')\,\rho^{e}(x)}\notag\\
&\quad{\color{blue}+2\;\mathcal{K}(x-w)\cdot\mathcal{K}(x'-w')\,\tfrac{1}{(x-z)^{2}}\;f^{abc}\,f^{cde}\,U^{ah}(z)\,U^{bl}(x)\,U^{dh}(x)\,U^{ml}(x')\,\rho^{m}(x')\,\rho^{e}(x)}\notag\\
&\quad{\color{blue}-2\;\mathcal{K}(x-w)\cdot\mathcal{K}(x'-w')\,\tfrac{1}{(x-z)^{2}}\;f^{abc}\,f^{cde}\,U^{hd}(x)\,U^{hl}(w')\,U^{ma}(z)\,U^{mb}(x)\,\rho^{l}(x')\,\rho^{e}(x)}\notag\\
&\quad{\color{blue}-2\;\mathcal{K}(x-w)\cdot\mathcal{K}(x'-w')\,\tfrac{1}{(x-z)^{2}}\;f^{abc}\,f^{cde}\,U^{hd}(x)\,U^{hl}(x')\,U^{ma}(z)\,U^{mb}(x)\,\rho^{l}(x')\,\rho^{e}(x)}\notag\\
&\quad{\color{blue}+\;\mathcal{K}(x-w)\cdot\mathcal{K}(x'-w')\,\tfrac{1}{(x-z)^{2}}\;f^{abc}\,f^{cde}\,U^{ha}(w)\,U^{hl}(x')\,U^{mb}(z)\,U^{md}(x)\,\rho^{l}(x')\,\rho^{e}(x)}\notag\\
&\quad{\color{blue}+\;\mathcal{K}(x-w)\cdot\mathcal{K}(x'-w')\,\tfrac{1}{(x-z)^{2}}\;f^{abc}\,f^{cde}\,U^{hb}(w)\,U^{hl}(x')\,U^{ma}(x)\,U^{md}(z)\,\rho^{l}(x')\,\rho^{e}(x)}\notag\\
&\quad{\color{blue}+\;\mathcal{K}(x-w)\cdot\mathcal{K}(x'-w')\,\tfrac{1}{(x-z)^{2}}\;f^{abc}\,f^{deh}\,U^{ad}(z)\,U^{bh}(x)\,U^{el}(x')\,\rho^{l}(x')\,\rho^{c}(x)}\notag\\
&\quad{\color{blue}+3\;\mathcal{K}(x-w)\cdot\mathcal{K}(x'-w')\,\tfrac{1}{(x-z)^{2}}\;f^{abc}\,f^{deh}\,U^{ah}(x)\,U^{dl}(w')\,U^{eb}(z)\,\rho^{l}(x')\,\rho^{c}(x)}\notag\\
&\quad{\color{blue}+2\,N_c\;\mathcal{K}(x-w)\cdot\mathcal{K}(x'-w')\,\tfrac{1}{(x'-z)^{2}}\;U^{ab}(x')\,U^{ac}(w)\,\rho^{b}(x')\,\rho^{c}(x)}\notag\\
&\quad{\color{blue}-2\,N_c\;\mathcal{K}(x-w)\cdot\mathcal{K}(x'-w')\,\tfrac{1}{(x'-z)^{2}}\;U^{ab}(x')\,U^{ac}(x)\,\rho^{b}(x')\,\rho^{c}(x)}\notag\\
&\quad{\color{blue}-\;\mathcal{K}(x-w)\cdot\mathcal{K}(x'-w')\,\tfrac{1}{(x'-z)^{2}}\;f^{abc}\,f^{cde}\,U^{hb}(x')\,U^{hl}(w)\,U^{ma}(z)\,U^{md}(x')\,\rho^{l}(x)\,\rho^{e}(x')}\notag\\
&\quad{\color{blue}+\;\mathcal{K}(x-w)\cdot\mathcal{K}(x'-w')\,\tfrac{1}{(x'-z)^{2}}\;f^{abc}\,f^{cde}\,U^{hb}(x')\,U^{hl}(x)\,U^{ma}(z)\,U^{md}(x')\,\rho^{l}(x)\,\rho^{e}(x')}\notag\\
&\quad{\color{blue}+\;\mathcal{K}(x-w)\cdot\mathcal{K}(x'-w')\,\tfrac{1}{(x'-z)^{2}}\;f^{abc}\,f^{cde}\,U^{hd}(x')\,U^{hl}(w)\,U^{ma}(z)\,U^{mb}(x')\,\rho^{l}(x)\,\rho^{e}(x')}\notag\\
&\quad{\color{blue}-\;\mathcal{K}(x-w)\cdot\mathcal{K}(x'-w')\,\tfrac{1}{(x'-z)^{2}}\;f^{abc}\,f^{cde}\,U^{hd}(x')\,U^{hl}(x)\,U^{ma}(z)\,U^{mb}(x')\,\rho^{l}(x)\,\rho^{e}(x')}\notag\\
&\quad{\color{blue}+\tfrac{1}{2}\;\mathcal{K}(x-w)\cdot\mathcal{K}(x'-w')\,\mathcal{K}(x-z)\cdot\mathcal{K}(x'-z)\;f^{abc}\,f^{dbe}\,U^{ha}(w')\,U^{hd}(w)\,\rho^{c}(x')\,\rho^{e}(x)}\notag\\
&\quad{\color{blue}+\tfrac{1}{2}\;\mathcal{K}(x-w)\cdot\mathcal{K}(x'-w')\,\mathcal{K}(x-z)\cdot\mathcal{K}(x'-z)\;f^{abc}\,f^{deh}\,U^{da}(w')\,U^{be}(z)\,U^{hl}(x)\,\rho^{c}(x')\,\rho^{l}(x)}\notag\\
&\quad{\color{blue}+\tfrac{1}{2}\;\mathcal{K}(x-w)\cdot\mathcal{K}(x'-w')\,\mathcal{K}(x-z)\cdot\mathcal{K}(x'-z)\;f^{abc}\,f^{deh}\,U^{ad}(w)\,U^{eb}(z)\,U^{lc}(x')\,\rho^{l}(x')\,\rho^{h}(x)}\notag\\
&\quad{\color{blue}+\tfrac{3}{2}\,N_c\;\mathcal{K}(x-w)\cdot\mathcal{K}(x'-w')\,\mathcal{K}(x-z)\cdot\mathcal{K}(x'-z)\;U^{ab}(x')\,U^{bc}(x)\,\rho^{a}(x')\,\rho^{c}(x)}\notag\\
&\quad{\color{blue}-\;\mathcal{K}(x-w)\cdot\mathcal{K}(x'-w')\,\mathcal{K}(x-z)\cdot\mathcal{K}(x'-z)\;f^{abc}\,f^{deh}\,U^{eb}(w)\,U^{lh}(x')\,U^{da}(x)\,\rho^{l}(x')\,\rho^{c}(x)}\notag\\
&\quad{\color{blue}-\;\mathcal{K}(x-w)\cdot\mathcal{K}(x'-w')\,\mathcal{K}(x-z)\cdot\mathcal{K}(x'-z)\;f^{abc}\,f^{ade}\,U^{hb}(w)\,U^{hd}(x')\,\rho^{e}(x')\,\rho^{c}(x)}\notag\\
&\quad{\color{blue}+2\;\mathcal{K}(x-w)\cdot\mathcal{K}(x'-w')\,\mathcal{K}(x-z)\cdot\mathcal{K}(x'-z)\;f^{abc}\,f^{ade}\,U^{hb}(x)\,U^{hd}(x')\,\rho^{e}(x')\,\rho^{c}(x)}\notag\\
&\quad{\color{blue}+\;\mathcal{K}(x-w)\cdot\mathcal{K}(x'-w')\,\mathcal{K}(x-z)\cdot\mathcal{K}(x'-z)\;f^{abc}\,f^{dae}\,U^{hb}(w')\,U^{hd}(x)\,\rho^{e}(x)\,\rho^{c}(x')}\notag\\
&\quad{\color{blue}-\;\mathcal{K}(x-w)\cdot\mathcal{K}(x'-w')\,\mathcal{K}(x-z)\cdot\mathcal{K}(x'-z)\;f^{abc}\,f^{deh}\,U^{eb}(x')\,U^{lh}(x)\,U^{da}(z)\,\rho^{l}(x)\,\rho^{c}(x')}\notag\\
&\quad{\color{blue}-\;\mathcal{K}(x-w)\cdot\mathcal{K}(x'-w')\,\mathcal{K}(x-z)\cdot\mathcal{K}(x'-z)\;f^{abc}\,f^{deh}\,U^{da}(z)\,U^{lc}(x')\,U^{be}(x)\,\rho^{h}(x)\,\rho^{l}(x')}\Bigg]\label{eq:four}
\end{align}
This list, read back by the parser, equals the checked \texttt{4rho.tex} lists ($1.2\times10^{-16}$ for $U=U^T$).

%==========================================================================================
\section{P2: which logarithm multiplies BFKL}\label{sec:p2}
Every row is an integral over the plus momentum $p^+$ of the unobserved gluon, with limits built from $\Lambda$, $k^+$ and $\vee$:
\begin{itemize}
\item the lower end, $p^+\to\Lambda$ (the unobserved gluon softer than the measured one), gives $\log\frac{k^+}{\Lambda}$; summed over all rows, its coefficient is JIMWLK$\otimes$LO, $\mathcal J$;
\item the upper end, $p^+\to\vee$ or $p^++k^+\to\vee$ (the unobserved gluon harder), gives $\log\frac{\vee}{k^+}$; summed over all rows, its coefficient is BFKL$\otimes$LO, $\mathcal B$ (earlier notes).
\end{itemize}
So the large logarithms of the NLO result are, exactly,
\begin{equation}
\frac{d^3N_{\rm LL}}{d^3k}=\NN\int_{x,x',z,w,w'}e^{-ik\cdot(w'-w)}\Big[\log\frac{k^+}{\Lambda}\,\mathcal J+\log\frac{\vee}{k^+}\,\mathcal B\Big].
\end{equation}
With your rule A, $\log\frac{k^+}{\Lambda}=\log\frac{\vee}{\Lambda}-\log\frac{\vee}{k^+}$:
\begin{equation}
\boxed{\frac{d^3N_{\rm LL}}{d^3k}=\NN\int_{x,x',z,w,w'}e^{-ik\cdot(w'-w)}\Big[\log\frac{\vee}{\Lambda}\,\mathcal J+\log\frac{\vee}{k^+}\,\big(\mathcal B-\mathcal J\big)\Big]}
\label{eq:total}
\end{equation}
\begin{itemize}
\item $\mathcal J$ is the bracket of your Sec.~3: the coefficient of $\log\frac{\vee}{\Lambda}$, i.e.\ the sum of all terms of Secs.~1, 2 (with P1).
\item $\mathcal B$ is the bracket of your Sec.~4: the coefficient of $\log\vee$, not of $\log\frac{\vee}{\Lambda}$.
\item At mid rapidity, $k^+=\sqrt{\vee\Lambda}$: $\frac12\log\frac{\vee}{\Lambda}\,(\mathcal J+\mathcal B)$.
\end{itemize}
In your Sec.~4 (all three equations): \old{$\log\Big[\frac{\vee}{\Lambda}\Big]$} $\to$ \new{$\log\Big[\frac{\vee}{k^+}\Big]$}, and the total \eqref{eq:total} is to be stated, otherwise Sec.~3 + Sec.~4 double count $\log\frac{\vee}{\Lambda}$.

%==========================================================================================
\section{P3: the prefactor of the singlet forms}\label{sec:p3}
With the projectile in a colour singlet, $\rho^b(x')\rho^c(x)\to\delta^{bc}\langle\rho^f(x)\rho^f(y)\rangle/(N_c^2-1)$ (then $x'\to y$), your Eq.~(3.4) reads
\begin{equation}
\frac{d^3N_{\rm NLO}}{d^3k}=\NN\log\Big[\frac{\vee}{\Lambda}\Big]\new{\frac{1}{N_c^2-1}}\int_{x,\old{x'}\new{y},z,w,w'}e^{-ik\cdot(w'-w)}\,\new{\langle\rho^f(x)\rho^f(y)\rangle}\,\Bigg[\mathcal K(x-w)\cdot\mathcal K(y-w')\,K(w,w',z)\Big[\dots\Big]+\dots\Bigg],
\end{equation}
and the same for Eq.~(4.3), with \new{$\log\big[\frac{\vee}{k^+}\big]$} (P2). Define once, before Eq.~(3.4):
\begin{equation}
\new{K(a,b,z)\equiv\mathcal K(a-z)\cdot\mathcal K(b-z)=\frac{(a-z)\cdot(b-z)}{(a-z)^2(b-z)^2},\qquad K(a,a,z)=\frac{1}{(a-z)^2}.}
\end{equation}
With this, $\NN\log\frac{\vee}{\Lambda}/(N_c^2-1)=P/k^+$ with $P=\frac{\alpha_s^2\delta Y}{8\pi^6(N_c^2-1)}$ of your App.~F.

%==========================================================================================
\section{P4: labels}\label{sec:p4}
\begin{itemize}
\item \textbf{Group I Row V.} The four lists are $d^3N_{2a}$, $d^3N_{2b\,\rm UV}$, $d^3N_{2b\,\rm finite}$ and a fourth one, also called $d^3N_{2b\,\rm finite}$.
\begin{itemize}
\item The fourth is $-d^3N_{2b\,\rm UV}$, term by term: \old{$d^3N_{2b\,\rm finite}$} $\to$ \new{$d^3N_{2b\,\rm CT}$} (counterterm).
\item The sum of all four is the checked Row V (Table~\ref{tab:gen}).
\end{itemize}
\item \textbf{``(with c.c.)''.}
\begin{itemize}
\item In Sec.~1 the heading of Group I no longer says ``(with c.c.)''; the heading of Group III does.
\item In both groups the c.c.\ is already written out in the lists, as it is in all Sec.~2 lists with this heading.
\item A note such as ``c.c.\ included'' would avoid adding it twice.
\end{itemize}
\end{itemize}

%==========================================================================================
\section*{Files}
In \texttt{nlo\_gluon\_paper/check\_v4/}:
\begin{itemize}
\item \texttt{parse\_v4.py}: parser of all terms and of the trace forms;
\item \texttt{check\_all.py}: Tables~\ref{tab:gen}--\ref{tab:s34};
\item \texttt{check\_singlet.py}: Table~\ref{tab:sing};
\item \texttt{check\_m1.py}, \texttt{m1\_singlet.py}: P1;
\item \texttt{make\_corrected.py}: the corrected copy, parsed back;
\item \texttt{make\_check.py}: this note.
\end{itemize}
\end{document}
